% A model section for the paper, written entirely in reduced form.
%
% Purpose: take a stance where the project has been mixing two stances. The
% structural system in claude/spread-explainer/spread_system.tex writes the
% spread as a function of latent objects (location premium, lease rates in two
% locations, a tariff hazard). Several of those are not identified with any
% data we can obtain. This note instead defines every object the paper uses as
% a projection of observable prices, and states what that costs us.
%
% No empirical work. The one numerical passage is an arithmetic illustration of
% the definitions, using quantities already computed elsewhere in the project,
% and is labelled as such.
%
% Build (from the repo root), keeping aux files out of the working tree:
%   pdflatex -interaction=nonstopmode -output-directory=<build> \
%     claude/spread-reduced-form/REDUCED_FORM_MODEL.tex        % run it twice
%
% Use pdflatex directly rather than latexmk: MiKTeX on this machine prints
% "You are running MiKTeX on an unsupported version of Windows" and exits
% non-zero even when the compile succeeded.
%
% Package set is the one the repo's other .tex files already compile with.

\documentclass[11pt,a4paper]{article}

\usepackage[margin=1in]{geometry}
\usepackage{amsmath}
\usepackage{amssymb}
\usepackage{url}
\usepackage{booktabs}
\usepackage{enumitem}
\usepackage{parskip}
\usepackage[hidelinks]{hyperref}

\newcommand{\phat}{\widehat{p}_t}
\newcommand{\omhat}{\widehat{\omega}_t}
\newcommand{\Var}{\operatorname{Var}}
\newcommand{\Cov}{\operatorname{Cov}}

\title{\textbf{The Reduced-Form Model}\\[0.4em]
\large From the Quoted Spread to a Measured New York Premium,\\
and from that Premium to Metal in Motion}
\author{Model section, GOLD project}
\date{22 September 2026}

\begin{document}
\maketitle

\noindent
This section builds the paper's price variable from the number a trader
quotes, shows why that number cannot be used as it stands, removes its time
dimension by estimating carry rather than assuming it, and arrives at a
premium that is defined by a projection of observed prices. It then connects
that premium to physical shipments through a threshold. Every object below is
a functional of observables. Nothing requires a lease rate, a storage
schedule, or a model of tariff risk.

\medskip
\noindent\textbf{The stance, stated once.} The project has two models in it and
has been sliding between them. One is \emph{structural}: the spread is written
as a function of latent objects --- a location premium, a lease rate in each
location, a hazard rate for a tariff --- and one asks which combinations the
data identify. The other is \emph{reduced form}: an object is defined by the
regression that produces it, and the paper's claims are claims about that
object's behaviour. The two are not interchangeable, and the current headline
variable is a hybrid of them, which is the worst of the three options. This
section goes fully reduced form and says in \S\ref{sec:stance} what that costs.

\tableofcontents

\section{The starting point: the quoted spread}
\label{sec:start}

Let $S_t$ be the London benchmark price of gold at date $t$, in US dollars per
troy ounce, for \emph{unallocated} metal in a London vault --- a claim on a
bank's pooled holding, the gold analogue of a deposit rather than a safe
deposit box. Let $F_t(\tau)$ be the price at the same date of a New York
futures contract that delivers metal $\tau$ days hence, where delivery means a
warrant: a title document to specific bars in an exchange-licensed New York
vault.

The number the market quotes is the difference:
\begin{equation}
b_t(\tau) \;=\; F_t(\tau) - S_t .
\end{equation}

The economic intuition the paper needs is already here. If New York metal is
dear relative to London metal by more than it costs to ship metal between
them, someone ships metal, and that shipment is recorded in trade statistics
as an export and an import even though nothing about world gold consumption
has changed. The paper's job is to make $b$ into a variable that can carry that
argument.

\section{Three reasons $b_t(\tau)$ cannot be used as it stands}
\label{sec:breaks}

\subsection*{It is denominated in a unit that is not comparable over time}

$b$ is dollars per ounce, and the gold price roughly quadrupled across the
sample. The same proportional gap between the two markets is about four times
as many dollars at the end of the period as at the start. A regression that
pools 2015 dollars with 2026 dollars is pooling incomparable units. This is
fixed by working in logs, or equivalently by scaling by $S_t$.

\subsection*{It contains carry, and carry moves with the delivery calendar}

$F$ is a price for delivery later. Even if the two markets are perfectly
integrated and metal is worth exactly the same in both places, $F$ exceeds $S$
by the cost of carrying metal from now until delivery, and that cost is
proportional to $\tau$.

The difficulty is that $\tau$ is not a fixed feature of the data. Gold futures
deliver in February, April, June, August, October and December, so the
contract the market actually uses has a $\tau$ that falls from roughly two
months to zero and then jumps back up when traders \emph{roll} --- close the
expiring contract and open the next one. Carry inherits that sawtooth. At a
4.5\% annual holding cost and a \$4{,}400 gold price, carry is about \$33 an
ounce at sixty days and under \$3 at five days. That swing is produced by the
delivery calendar and nothing else, and it is the same order of magnitude as
the market dislocations the paper is trying to detect.

\subsection*{It mixes a level with a rate}

Two economically different things sit inside $b$.

\begin{itemize}[leftmargin=1.2em,itemsep=2pt,topsep=2pt]
\item A \textbf{level}: whatever makes metal in New York worth more than metal
in London at a given moment --- location, the bar form each market trades,
the difference between an exchange warrant and a claim on a bank. A level is
paid once. It does not care whether the contract expires in ten days or a
hundred.
\item A \textbf{rate}: financing, storage, and the income forgone by not
lending the metal out. A rate accrues per unit of time and therefore scales
with $\tau$.
\end{itemize}

Any single-number summary of $b$ must choose a horizon at which to compare the
two, and that choice determines the answer. The project has already seen this:
the same day reads as $6.9$, $2.8$ or $1.8$ percentage points annualised
depending on whether the curve is read at thirty, ninety or one hundred and
eighty days.

\medskip
\noindent The conclusion is not that $b$ is the wrong idea. It is that $b_t(\cdot)$ is a
\emph{function of $\tau$}, not a scalar, and the paper's variable has to be a
feature of that function rather than one of its values.

\section{Carry, and the decision to estimate it}
\label{sec:carry}

\subsection*{What carry is}

Carrying an ounce from today to date $\tau$ costs the interest on the money
tied up, plus the cost of storing and insuring the metal, less the income the
owner forgoes by not lending the metal out. That last term is the \emph{lease
rate}: metal, unlike most commodities, can be lent, and the borrower pays for
it.

For a macroeconomist the whole structure is covered interest parity with gold
in the role of the foreign currency. $S$ is the spot exchange rate, $F$ the
forward rate, the dollar interest rate is the domestic rate, and the lease
rate is the foreign rate. Carry is the forward points. Deviations from it are
the gold analogue of a cross-currency basis.

\subsection*{Why the conventional fix is not reduced form}

The usual repair is to convert the spread into the annual financing rate it
implies and then subtract an assumed carry:
\begin{align}
\text{implied rate}_t(\tau) &= \left(\frac{F_t(\tau)}{S_t}-1\right)\frac{365}{\tau},\\
\text{dislocation}_t(\tau) &= \text{implied rate}_t(\tau) - \bigl(r_t + \bar{s}\bigr).
\end{align}

This has two defects, and naming them is the point of this section.

\textbf{First, it is a hybrid.} The implied rate is a reduced-form measurement:
it is a transformation of two observed prices. The subtracted term is a
structural input: it asserts a storage cost $\bar s$ that is not observed and a
lease rate of zero that is certainly wrong. So the headline variable is a
measurement corrected by assumptions, and it inherits the fragility of the
assumptions without gaining any of the discipline that a structural model
would provide in exchange. Both errors are also signed and trending: storage is
billed in dollars per ounce, so as a \emph{rate} it falls when the gold price
rises, and a constant $\bar s$ therefore overstates carry by a growing amount
across a sample in which gold quadrupled.

\textbf{Second, it annualises a level.} Multiplying by $365/\tau$ is the right
operation for something that accrues over time and the wrong one for something
paid once. A constant level premium, put through this transformation, produces
a fresh sawtooth pointing the other way, and diverges as $\tau\to 0$. The
repair cures one calendar artefact by manufacturing another.

\subsection*{The reduced-form alternative}

Do not assume carry. \emph{Estimate} it. On any date the market quotes prices
for several delivery months at once, and those prices trace out a term
structure: a yield curve for New York gold. The slope of that curve is the
market's own carry, revealed by the prices themselves. Nothing has to be
assumed about storage, and no interest-rate series is needed to extract it.

\section{The estimator}
\label{sec:estimator}

\subsection*{The daily cross-section}

On date $t$, collect every listed contract $i=1,\dots,n_t$ with a settlement
price $F_{it}$, a horizon $\tau_{it}$ and open interest $\mathrm{OI}_{it}$
(the number of contracts outstanding, a liquidity measure). Run the weighted
least-squares regression
\begin{equation}
\ln F_{it} \;=\; a_t \;+\; b_t\,\tau_{it} \;+\; u_{it},
\qquad w_{it}=\sqrt{\mathrm{OI}_{it}} .
\label{eq:cs}
\end{equation}

The functional form follows from compounding, not from a behavioural
assumption: if metal is carried at a constant proportional cost $c$ per day
then $F(\tau)=F(0)e^{c\tau}$, so $\ln F$ is linear in $\tau$ with slope $c$ and
intercept $\ln F(0)$. The intercept is the extrapolation of the New York
forward curve back to zero horizon --- the price today of New York-deliverable
metal.

Define the paper's two price variables:
\begin{align}
\phat &\;\equiv\; a_t - \ln S_t
  &&\text{the \emph{New York premium}: a level, in log points of spot,}\label{eq:phat}\\
\omhat &\;\equiv\; 365\,b_t - r_t
  &&\text{\emph{excess carry}: a rate, in points a year.}\label{eq:omhat}
\end{align}

\subsection*{What this construction achieves}

\begin{enumerate}[leftmargin=1.4em,itemsep=3pt,topsep=2pt]
\item \textbf{No rate data enters $\phat$.} Carry loads entirely on $\tau$ and
is absorbed by $b_t$. Whatever the true financing cost, storage fee or lease
rate, they shift the slope, not the intercept. The premium is therefore free of
every unobservable that has been blocking the project. Only $\omhat$ needs an
interest rate, and only because subtracting $r_t$ is what makes a slope
interpretable.
\item \textbf{The roll cannot contaminate it.} Nothing here picks a front
month. Contracts enter as a cross-section, so the sawtooth has no route in, and
the short-$\tau$ filters that version~1 needed are no longer load-bearing.
\item \textbf{The horizon choice disappears.} We do not read the curve at a
chosen maturity; we report its intercept and its slope. The level and the rate
are reported as the different objects they are, which removes the arbitrariness
documented in \S\ref{sec:breaks}.
\item \textbf{$\phat$ is defined by the projection.} This is the crucial point
of the whole section. $\phat$ is not a parameter of a pricing model that may or
may not be true; it is a number computed from prices by a stated rule. It
exists, and means the same thing, whether or not any structural account of the
gold market is correct. A structural model is then a \emph{claim about what
$\phat$ should equal} --- which is a useful thing to have, and a different
thing from the estimand.
\end{enumerate}

\subsection*{Choices that have to be stated, because they are choices}

\begin{itemize}[leftmargin=1.2em,itemsep=2pt,topsep=2pt]
\item \emph{Which contracts.} Restrict to $15\le\tau\le 400$ days and open
interest above a floor. Contracts inside fifteen days are in the delivery
period, where price behaviour is about warrants rather than carry; far-dated
and thinly held contracts carry settlement prices that are marked rather than
traded. This leaves a median of about six contracts a day.
\item \emph{Weights.} Thin contracts have noisier settlements, so weighting by
$\sqrt{\mathrm{OI}}$ is a proxy for inverse standard deviation. It is an
assumption about the error structure, not a result, and the estimates should be
shown unweighted as well.
\item \emph{Linear, not quadratic.} A horizon-dependent risk --- the chance a
tariff lands before delivery --- would curve the term structure, and fitting
$\tau^2$ is tempting. With six points a day a quadratic is fragile and its
coefficient is close to collinear with the slope. Baseline linear; curvature as
a robustness exercise and as a diagnostic, never as a headline.
\item \emph{Daily.} Equation~\eqref{eq:cs} is estimated separately each day
rather than pooled with time fixed effects, because both $a_t$ and $b_t$ are
the objects of interest and neither is a nuisance.
\end{itemize}

\subsection*{A sharper measurement of the same estimand}

If a forward curve for \emph{London} metal is available at matched dates, the
same projection can be run on it:
$\ln F^{L}_{it} = a^{L}_t + b^{L}_t \tau_{it} + u^{L}_{it}$. Then
\begin{equation}
\widehat{\pi}_t \;\equiv\; a_t - a^{L}_t
\label{eq:twocurve}
\end{equation}
is the log premium of New York metal over London metal at zero horizon,
measured entirely inside futures markets. It needs no spot benchmark, no
interest rate, and no storage figure; and if the two curves settle at the same
instant it is free of the timing problem in \S\ref{sec:timing}. The slope
difference $b_t-b^{L}_t$ is the cross-location carry differential, which is the
natural place to look for a squeeze in the London lending market.

This is not a different model. It is equation~\eqref{eq:phat} measured better,
and it should be presented that way: the baseline result uses $\phat$, and
$\widehat{\pi}_t$ is the robustness check that removes three assumptions at
once.

\section{What the premium contains, and why we do not decompose it}
\label{sec:contains}

$\phat$ contains, at least: the value of metal being in New York rather than
London when policy might restrict moving it; the difference between the bar
forms the two markets trade; the difference between an exchange warrant and an
unallocated claim on a bank; and measurement error from pricing the two legs at
different moments.

The reduced-form stance is that we do not attempt to separate these, and the
reason is not modesty but identification: with price data alone, a constant
institutional wedge and a constant risk premium are observationally
equivalent. Attributing a share of $\phat$ to any one component would be
asserting, not estimating.

What makes this stance workable is that the paper's three claims are invariant
to the decomposition:

\begin{enumerate}[leftmargin=1.4em,itemsep=2pt,topsep=2pt]
\item $\phat$ \emph{moves on dated policy news}, which establishes that it is
not a constant institutional wedge;
\item $\phat$ \emph{predicts physical shipments}, and only above a threshold,
which establishes that it is the price signal that actually moves metal;
\item the shipments it predicts \emph{are relocation rather than absorption},
which is what makes them a contaminant of trade statistics.
\end{enumerate}

None of the three requires knowing how much of the premium is bar form and how
much is tariff risk. What we give up is equally clear and should be written in
the paper rather than left for a referee to notice: we cannot decompose the
premium, so we make no statement about the price of tariff risk as such, and no
welfare claim that would require one.

\section{Measurement error, and the one structural fact we do use}
\label{sec:timing}

The two legs are not observed at the same instant. The London benchmark is
struck in the afternoon London auction; the New York settlement is calculated
about three and a half hours later. Let $\varepsilon_t$ be gold's return over
that window.

One fact about the trading mechanism does real work here, and it is worth
keeping precisely because it is a fact about the mechanism rather than an
assumption about behaviour: \emph{every contract settles at the same instant}.
So $\varepsilon_t$ multiplies all of that day's futures prices equally. In
logs,
\begin{equation}
\widehat{a}_t = a_t + \varepsilon_t,
\qquad
\widehat{b}_t = b_t .
\end{equation}
The timing error is a pure \emph{level} contaminant. It biases the premium by
$\mathbb{E}[\varepsilon_t]$ and adds $\Var(\varepsilon_t)$ to its variance, and
it leaves the estimated carry slope untouched.

Two consequences follow, and they determine the frequency at which each part of
the paper should be run.

\begin{itemize}[leftmargin=1.2em,itemsep=3pt,topsep=2pt]
\item \textbf{As an outcome}, in the news regression, classical measurement
error inflates standard errors but does not bias the coefficient --- provided
the news does not arrive inside the three-and-a-half-hour window, which for
some announcements it does. The defence is an event window that opens before
and closes after the affected settlement, plus re-timing the New York leg to
the London auction instant where intraday data allow.
\item \textbf{As a regressor}, in the flow equation, the same error attenuates
the slope towards zero and, because the specification has a kink, smooths the
kink and biases the estimated threshold. The defences are averaging to monthly
frequency, which shrinks the noise variance by the number of days in the month;
re-timing; and instrumenting the premium with dated announcements.
\end{itemize}

That the news regression belongs at daily frequency and the flow regression at
monthly is therefore not a convenience. It falls out of where the measurement
error lands.

\section{From premium to metal: the second reduced form}
\label{sec:flow}

Write the premium in dollars per ounce, $P_t = \phat\,S_t$, and let $Q_t$ be
tonnes shipped along the relevant corridor in month $t$. Metal moves only when
the premium covers the all-in cost of moving it: freight, insurance, recasting
into the bar form the destination market accepts, and the financing of metal in
transit. Below that cost, nothing happens; above it, the trade is profitable
and scales with how far above it sits. The reduced-form representation is a
hinge:
\begin{equation}
Q_t \;=\;
\beta^{W}\max\!\left(0,\;P_t-\kappa^{W}\right)
\;-\;\beta^{E}\max\!\left(0,\;-P_t-\kappa^{E}\right)
\;+\;\gamma' x_t \;+\; u_t .
\label{eq:hinge}
\end{equation}

Four features of this deserve a sentence each.

\textbf{Dollars, not rates.} Freight, insurance and refining are charged by
weight, so the threshold is naturally a dollar amount per ounce and should not
be expressed as a percentage. This is a prior about the technology, not a
fitted result, and the paper should say so plainly; the two specifications are
only distinguishable when the gold price moves a lot.

\textbf{A kink, not a slope.} A linear specification lets the many months in
which nothing moves pull down the estimated response of the months in which a
great deal moves. The kink is also the economically interesting parameter:
$\kappa$ is the transatlantic transfer cost, which is not published anywhere,
so estimating it is a contribution in its own right rather than a nuisance.

\textbf{Two sides, separately.} Above the band New York is dear and metal flows
west; below it New York is cheap and metal flows east. Both regimes occur. A
single dummy or a symmetric specification would impose that the two thresholds
and the two slopes are equal, which is exactly the kind of restriction the
data can test.

\textbf{Daily hinge, then aggregate.} Because $\max(0,\cdot)$ is convex,
applying the hinge to a monthly average premium is not the same as averaging
the daily hinge, and it systematically understates the response: a month
containing a few days far above the threshold and many days below it looks, in
its average, like a month that never crossed. Apply the hinge day by day and
aggregate afterwards.

Two further shifters belong in $x_t$, and both are reduced-form controls rather
than structural amendments. Capacity --- refineries and secure freight can only
move so much per month --- censors $Q$ from above, so a period of sustained
high premia with flat flows is not evidence against the mechanism. And the
stock of metal already inside the destination market matters: an arbitrage can
be closed by re-warranting metal that is already in New York, which produces no
cross-border shipment at all. That is not a defect of the specification but the
central asymmetry of the whole project, and it belongs in the flow equation as a
control and in the accounting as a bound.

\section{How the pieces fit together}
\label{sec:chain}

\begin{center}
\begin{tabular}{@{}p{0.055\textwidth}p{0.30\textwidth}p{0.20\textwidth}p{0.33\textwidth}@{}}
\toprule
\textbf{Step} & \textbf{What is estimated} & \textbf{Frequency} & \textbf{What it establishes} \\
\midrule
1 & Cross-sectional projection~\eqref{eq:cs}; the premium~\eqref{eq:phat} and
excess carry~\eqref{eq:omhat} & daily & a price variable free of carry, the
roll and any horizon choice \\
2 & Response of $\Delta\phat$ to signed, timed policy announcements & daily &
the premium is priced policy risk, not a constant wedge \\
3 & The hinge~\eqref{eq:hinge}: $\beta$ and $\kappa$ & monthly & the premium
moves metal, and only once it clears a cost \\
4 & Fitted flows against absorption benchmarks and vault stocks & monthly &
those flows are relocation, not consumption \\
5 & Relocation flows applied to reported trade aggregates & monthly & the size
of the contamination --- the paper's payoff \\
\bottomrule
\end{tabular}
\end{center}

Steps 2 and 3 can be run jointly, with the announcements instrumenting the
premium in the flow equation. That buys consistency if the premium is measured
with error, at the cost of an exclusion restriction: announcements must affect
shipments only through the premium. Three things threaten it, and each has a
reduced-form remedy rather than a structural one --- safe-haven demand shifting
flows directly, which is why ETF and central-bank series enter as controls;
precautionary relocation by owners who would move metal whatever the price,
which is untestable and must be conceded; and country-wide tariff measures that
affect all trade with a partner, which is why gold-specific events are
preferred to broad ones.

\section{Arithmetic illustration of the definitions}
\label{sec:illustration}

No estimation is reported in this note. To make the mapping concrete, here is
one date run through the definitions, using quantities computed earlier in the
project.

On 29 January 2025 the fitted curve had an intercept implying a zero-horizon
New York price of \$2{,}769.38 against a London benchmark of \$2{,}756.30. Then
\[
\phat=\ln\!\left(\frac{2769.38}{2756.30}\right)=0.473\%,
\qquad
P_t=\phat\,S_t=\$13.08\text{ per ounce},
\]
which, against a transfer cost of order \$3 an ounce, sits well inside the
region where equation~\eqref{eq:hinge} predicts metal moves. The fitted slope on
the same day was $5.22\%$ a year, so with an overnight rate of $4.35\%$ the
excess carry was $\omhat\approx+0.9$ points before any storage adjustment.

Two scale conversions make these numbers readable. One per cent of spot at a
\$4{,}500 gold price is \$45 an ounce, or \$1.45 million a tonne. A \$3 per
ounce threshold is about \$96{,}000 a tonne, which is where air freight,
insurance and recasting plausibly sit --- a sanity check on the estimand, not a
validation of it.

\section{Notation, units and signs}
\label{sec:units}

\begin{center}
\begin{tabular}{@{}lll@{}}
\toprule
\textbf{Symbol} & \textbf{Meaning} & \textbf{Unit} \\
\midrule
$S_t$ & London benchmark, unallocated metal & \$/oz \\
$F_t(\tau)$ & New York futures price, delivery in $\tau$ days & \$/oz \\
$\tau$ & calendar days to the delivery-relevant date & days \\
$a_t,\ b_t$ & intercept and slope of the daily curve~\eqref{eq:cs} & log points; log points/day \\
$\phat$ & New York premium, $a_t-\ln S_t$ & log points of spot \\
$P_t$ & the same premium in money, $\phat S_t$ & \$/oz \\
$\omhat$ & excess carry, $365b_t-r_t$ & points a year \\
$\kappa$ & transfer cost threshold & \$/oz \\
$\beta$ & flow response above the threshold & tonnes per month per \$/oz \\
$Q_t$ & metal shipped & tonnes a month \\
\bottomrule
\end{tabular}
\end{center}

\noindent\textbf{Conventions.} Rates are annualised decimals and differences
between rates are quoted in percentage points. Day count is 365 throughout; if
an overnight rate quoted on a 360-day basis is used in~\eqref{eq:omhat} it must
be scaled by $365/360$ first. Mass is tonnes, with $1$~t${}=32{,}150.7$~oz, and
prices are converted at the boundary rather than inside the analysis.

\noindent\textbf{Signs.} $\phat>0$ means New York is dear, so metal flows west,
towards New York. $\phat<0$ means the reverse. Both regimes occur in the
sample, and~\eqref{eq:hinge} is written to handle both rather than assuming the
westward case.

\section{Assumptions, and where each one fails}
\label{sec:assumptions}

\begin{center}
\begin{tabular}{@{}p{0.30\textwidth}p{0.30\textwidth}p{0.33\textwidth}@{}}
\toprule
\textbf{Assumption} & \textbf{Where it enters} & \textbf{How it fails} \\
\midrule
The log forward curve is linear in the horizon & \eqref{eq:cs} & horizon-dependent
risk curves it; a tariff hazard bends it downward, loading onto the slope and
biasing $\omhat$, not $\phat$ \\
The timing shock is common to all contracts & \S\ref{sec:timing} & fails if
contracts settle at different instants, or if a far month is marked rather than
traded; then the error tilts the curve and contaminates $\phat$ as well \\
Settlement noise falls with open interest & weights in~\eqref{eq:cs} & fails
when a thin far contract is pinned to a model price and so looks \emph{less}
noisy than it is \\
The threshold is constant in dollars within the estimation window &
\eqref{eq:hinge} & freight and insurance move with fuel and with wartime risk
premia; a level shift in $\kappa$ mimics a change in $\beta$ \\
Shipments are recorded in the month the premium was earned & aggregation &
exporters record on departure and importers on arrival, so month-boundary
crossings misalign the two sides \\
Announcements move flows only through the premium & step 3 of
\S\ref{sec:chain} & safe-haven demand and precautionary relocation both violate
it; the second is untestable with available data \\
The premium is a price signal, not a quantity constraint & \eqref{eq:hinge} &
when capacity binds, $Q$ stops responding and $\beta$ is attenuated; treat the
ceiling explicitly \\
\bottomrule
\end{tabular}
\end{center}

\section{What we give up by not going structural}
\label{sec:stance}

A structural version of this exercise would write the premium as a function of
a location risk premium, a lease rate in each market and a hazard rate for
tariff imposition, and would estimate those objects. It is worth being explicit
about why the paper does not do that.

\begin{enumerate}[leftmargin=1.4em,itemsep=3pt,topsep=2pt]
\item \textbf{The central parameter has no data behind it.} There is no
reported New York lending market for gold at any price. In the structural
system the New York lease rate appears only in combination with the storage
cost, so the futures curve alone can never separate them. A structural estimate
would be an assumption wearing an estimate's clothes.
\item \textbf{The components of the premium are observationally equivalent in
price data.} Location risk, bar-form scarcity and warrant-versus-claim trust all
enter as a level. No amount of price data separates them; only owner-level or
transaction-level data would, and it does not exist publicly.
\item \textbf{The claims the paper actually makes do not need the
decomposition.} Everything in \S\ref{sec:contains} survives any allocation of
the premium across its components.
\end{enumerate}

So the structural system belongs in an appendix, where it does two honest jobs:
it explains what $\phat$ and $\omhat$ would equal under an explicit pricing
model, which is how the reader knows what sign to expect; and it states
precisely which combinations are not identified, which is how the reader knows
the main text is not quietly claiming them.

\medskip
\noindent\textbf{Three consequences for the paper's prose, which are the
practical content of taking this stance.}

\begin{itemize}[leftmargin=1.2em,itemsep=3pt,topsep=2pt]
\item \emph{Retire ``dislocation'' as the headline variable.} It is the hybrid
object of \S\ref{sec:carry}: a measurement minus an assumption. Report $\phat$
and $\omhat$ instead, which are both projections of observables. If a single
annualised number is wanted for comparability with market commentary, report it
in a footnote with the assumed storage rate stated next to it.
\item \emph{Never attribute a share of $\phat$ to a named component.}
Sentences of the form ``the tariff premium reached $x$ per cent'' should read
``the New York premium reached $x$ per cent'', with the tariff attribution
carried by the event study rather than by the label.
\item \emph{Report the level and the rate separately, always.} They answer
different questions --- how much more a dealer will pay to have metal in New
York today, versus how expensive it is to hold it there --- and collapsing them
into one number was the source of the confusion this section exists to remove.
\end{itemize}

\end{document}
